% !PRINTED 200
\weberpage{200}
This substitution must occur in $G$, and by Theorem 1 we must have $c_1 = c_2$, and hence $\Theta_1$ identical with $\Theta_2$. The numbers $1, \varepsilon_1, \varepsilon_2, \ldots, \varepsilon_{\nu-1}$ must therefore together contain all the $\nu$th roots of unity, and hence coincide with the powers of a primitive one among them. Since, on the other hand, all the powers of one of the quantities $\Theta$ must occur in group (1), it follows that, if the $\varepsilon$ occurring in $\Theta$ is a primitive $\nu$th root of unity, the whole group (1) is exhausted by the powers of $\Theta$. % [sic/reference: both occurrences of ``group (1)'' are printed; the stabilizer group Q in (2), named in the next sentence, is required by the argument.]
Hence this group $Q$ consists of the elements

\weq{4}{1, \Theta, \Theta^2, \ldots, \Theta^{\nu-1}.}

If we take (following Theorem 2) $\Theta$ as a multiplicative substitution, then 0 is the second pole of $\Theta$, and at the same time the second pole of all the substitutions (4). This pole is therefore likewise at least \gesperrt{$\nu$-fold}. But since the two poles of $\Theta$ can be interchanged in this consideration, we obtain the theorems:

3. \gesperrt{A $\nu$-fold pole determines a cyclic group $Q$ of degree $\nu$ contained in $G$.}

4. \gesperrt{The two poles of any substitution $\Theta$ of the group $G$ are poles of equal multiplicity, and are the only two poles of the group $Q$.}

By \S.\,2, $\mu - 1$ elements

\weq{5}{\psi_1, \psi_2, \ldots, \psi_{\mu-1}}

can be so selected in $G$ that $\psi_1 Q, \psi_2 Q, \ldots, \psi_{\mu-1} Q$ are the cosets of $Q$, and hence

\weq{6}{G = Q + \psi_1 Q + \psi_2 Q + \cdots + \psi_{\mu-1} Q}

holds. We now set

\weq{7}{a_1 = \psi_1(a), a_2 = \psi_2(a), \ldots, a_{\mu-1} = \psi_{\mu-1}(a).}

The quantities $a_1, a_2, \ldots, a_{\mu-1}$ are not only all different from $a$, but also different from one another. For if, say,

\[ \psi_1(a) = \psi_2(a), \]

it would follow that

\[ a = \psi_1^{-1}\psi_1(a) = \psi_1^{-1}\psi_2(a), \]

and $\psi_1^{-1}\psi_2$ would be contained in $Q$, hence $\psi_2$ in $\psi_1 Q$, which is contrary to the assumption. It is now easy to show that these values $a_1, a_2, \ldots, a_{\mu-1}$ are all $\nu$-fold poles of the group. It suffices to show this for $a_1$.
